A battery only discharges under any admissible input.
The set of sufficient energy states therefore cannot be held invariant by control alone, and sufficiency is meaningful only relative to reaching a charging station.
Persistification frameworks certify this property with a CBF~\cite{notomista2018persistification, notomista2020persistification}.
Each robot's state is augmented with its battery level, and the controller keeps the robot inside the set from which a charger remains reachable on the remaining charge.
Extensions coordinate shared stations across a team~\cite{fouad2022energy}.
In all of these, the certificate serves one robot and one purpose, and its margin is discarded once sufficiency is established.
\textit{We reinterpret this margin as the energy a robot can spare and elevate it to the precedence cue that coordinates the team.}

The barriers of Section~\ref{sec:2.A} certify specifications that a controller can actively maintain. 
Energy sufficiency is not such a specification. 
Under any admissible input the battery strictly discharges, so the set $\{e\geq e_{\min}\}$ admits no control that holds it invariant, and the presumption behind Eq.~\eqref{eq:cbf_condition} fails.
Sufficiency acquires meaning only in relation to a charging station: a robot is sufficiently charged not when its level is high, but when what it retains still covers the cost of reaching a station.

Persistification frameworks make this precise by certifying reachability rather than level~\cite{notomista2018persistification, notomista2020persistification}.
The state is augmented with the battery level, and the invariant set is not $\{e\geq e_{\min}\}$ but the set of augmented states from which a station remains reachable on the charge in hand. 
A robot approaching depletion is turned back before its remaining charge falls below the return cost, and this turning-back is exactly a control that keeps the augmented state inside the reachable set. 
Forward invariance is thereby restored on the correct set, and a CBF certifies it in the standard form. 
Subsequent work extends the certificate to teams sharing a limited set of stations~\cite{fouad2022energy}.

Across these frameworks the energy certificate is single-purpose.
Its margin, the charge a robot holds beyond the return cost, is read only to confirm that sufficiency is met, and is discarded once the constraint is satisfied.
We retain this margin and read it as a measure of how much a robot can spare. 
A robot with a large margin can absorb a detour that a robot near its return cost cannot, so the margin ranks two robots by their capacity to yield. 
This reinterpretation turns a per-robot feasibility certificate into the pairwise signal that Section~\ref{sec:2.B} left unspecified: an allocation that is neither learned from demonstrations nor assigned by a designer, but computed onboard from the energy state itself.